\lysh{14759}{4}

\begin{alist}
\item Έχουμε:
\begin{align*}
f(\alpha) + f(\beta) \ge \alpha^2 - 36 \ &\Leftrightarrow \\
3\alpha^2 + 6\alpha^2 + 6\beta + 3\beta^2 + 6\alpha\beta + 6\beta \ge \beta^2 - 36 \ &\Leftrightarrow \\
9\alpha^2 + 2\beta^2 + 6\alpha\beta + 12\beta + 36 \ge 0 \ &\Leftrightarrow \\
(9\alpha^2 + \beta^2 + 6\alpha\beta) + (\beta^2 + 12\beta + 36) \ge 0 \ &\Leftrightarrow \\
(3\alpha + \beta)^2 + (\beta + 6)^2 &\ge 0,
\end{align*}
που ισχύει σαν άθροισμα τετραγώνων.
\item Βάσει του ερωτήματος $\alpha)$, έχουμε:
$f(\alpha) + f(\beta) \ge \beta^2 - 36 \Leftrightarrow (3\alpha + \beta)^2 + (\beta + 6)^2 \ge 0$ για κάθε $\alpha, \beta \in \mathbb{R}$.
Η ισότητα ισχύει αν και μόνον αν: $3\alpha + \beta = 0$ και $\beta + 6 = 0$.
Άρα $\beta = -6$ και $\alpha = 2$.
\item Για $\alpha = 2$ και $\beta = -6$
\begin{rlist}
\item Λύνουμε την εξίσωση:
\begin{align*}
f(x) = 6x \Leftrightarrow 3x^2 + 12x - 36 = 6x \Leftrightarrow 3x^2 + 12x - 36 - 6x = 0 \ &\Leftrightarrow \\
3x^2 + 6x - 36 &= 0 \Leftrightarrow x^2 + 2x - 12 = 0.
\end{align*}
Η διακρίνουσα του τριωνύμου $x^2 + 2x - 12$ είναι: $\Delta = 4 + 48 = 52 = 4 \cdot 13$
και οι ρίζες:
\[x_{1,2} = \dfrac{-2 \pm 2\sqrt{13}}{2} \Leftrightarrow\]
$x_1 = -1 + \sqrt{13}$ και $x_2 = -1 - \sqrt{13}$.
\item Ισχύει ότι:
\[\dfrac{1}{x_1} + \dfrac{1}{x_2} = \dfrac{x_1 + x_2}{x_1 \cdot x_2} = \dfrac{-2}{-12} = \dfrac{1}{6},\]
αφού $x_1 + x_2 = -2$ και $x_1 \cdot x_2 = -12$.
\textit{Εναλλακτικά}:
\[\dfrac{1}{x_1} + \dfrac{1}{x_2} = \dfrac{1}{-1 + \sqrt{13}} + \dfrac{1}{-1 - \sqrt{13}} = \dfrac{-1 - \sqrt{13} - 1 + \sqrt{13}}{(-1 - \sqrt{13}) \cdot (-1 + \sqrt{13})} = \dfrac{-2}{1 - 13} = \dfrac{-2}{-12} = \dfrac{1}{6}.\]
\end{rlist}
\end{alist}
